\begin{statement}{C.2.UN.45}{Proved}{Every compact curve is the ruling base of an actual normalization component}
For every smooth compact connected complex curve $C$, there is a smooth
compact complex threefold $Y$ and a proper map $f:Y\to\PP^1$ with connected
fibres, whose restriction over $\PP^1\setminus\{0\}$ is the product with a
fixed two-torus, and whose reduced zero fibre has Euler number zero and
contains a smooth multiplicity-one component
\[
                         R=\PP(N_{C/U})\longrightarrow C.
\tag{C.RU.1}
\]
Here $C\hookrightarrow U$ is the explicit vertical embedding constructed
below. The component $R$ is retained unchanged by the Euler-balancing
operations. Its birational class is that of $C\times\PP^1$.
Consequently every ruled surface birational class occurs in an actual
compact zero-Euler torus fibre. A rational component is also obtained
unchanged as the exceptional plane of a point blowup followed by balancing
away from that plane.

Writing $g=g(C)$, the integral cohomology of $R$ is torsion-free and has ranks
\[
              (b_0,b_1,b_2,b_3,b_4)(R)=(1,2g,2,2g,1),
              \qquad h^\bullet(R,\OO_R)=(1,g,0).
\tag{C.RU.2}
\]
The construction proves occurrence of every ruled birational class; it
does not prescribe an arbitrary rank-two bundle as the normal bundle in
(C.RU.1).
\end{statement}

\begin{proof}
Choose a nonconstant meromorphic function $x$ on $C$, supplied by
Riemann--Roch. It gives a finite map $C\to\PP^1$, so the meromorphic
function field of $C$ is a finite separable extension of $\C(x)$.
Choose a primitive element $y$. The meromorphic map
\[
                    C\longrightarrow\PP^2,\qquad [1:x:y]
\tag{C.RU.3}
\]
extends at every point: multiply its three local meromorphic coordinates
by a power of a uniformizer so that their minimum vanishing order is zero.
At least one resulting holomorphic coordinate is nonzero. The image is an
integral projective plane curve $D$, and the map is finite and birational,
because $\C(x,y)=\mathcal M(C)$. Thus $C$ is the actual normalization of
$D$. For $C=\PP^1$ one may choose $x=s,y=s^2$, giving a smooth conic.
In every case choose projective coordinates so that $D$ is not the fixed
line used below.

Start with $A\times\PP^1$, where $A$ is any fixed smooth complex two-torus,
and blow up a point of $A\times\{0\}$. Its exceptional divisor is an
actual $\PP^2$. Place $D$ in this plane. Successively blow up singular
points of its strict transform, performing each point blowup in the
ambient smooth threefold. The strict transform of the plane is precisely
its surface point blowup: in local coordinates this is the strict
transform of the coordinate plane in the three ordinary blowup charts.
The full map to $\PP^1$ is unchanged over its punctured base.

Here the point-resolution process is finite without an unproved
termination assertion. For an integral curve $D_0$ on a smooth surface,
normalization gives
\[
  p_a(D_0)-g(C)=\operatorname{length}
                (\nu_*\OO_C/\OO_{D_0})\ge0.
\tag{C.RU.4}
\]
At a singular point the multiplicity $m$ is at least two. Under a surface
point blowup, $D_1=\pi^*D_0-mE$ and
$K_{V_1}=\pi^*K_{V_0}+E$. Adjunction gives
\[
                        p_a(D_1)=p_a(D_0)-\frac{m(m-1)}2.
\tag{C.RU.5}
\]
The normalization is unchanged through this proper birational map.
Therefore its total nonnegative normalization length decreases by at
least one at every singular-point blowup. After finitely many such
blowups the strict transform is smooth and is the specified $C$ itself.
This gives a genuine closed embedding $C\hookrightarrow U$ into the
resulting smooth compact threefold.

The strict transform of the original exceptional plane has full-fibre
multiplicity one, since that was its coefficient after the first point
blowup and later strict transforms retain the coefficient. The curve
$C$ is not any boundary curve of this plane: it dominates the original
$D$, which is different from the line of intersection with the original
torus component, whereas every subsequently exceptional boundary curve
maps to a point of that plane. Consequently a generic point of $C$ lies
only in this multiplicity-one component. Blow up the smooth curve $C$ in
$U$. The exceptional divisor is exactly the projectivized normal bundle
in (C.RU.1), with coefficient one in the full fibre. This follows also
from the local charts of the blowup of $(x,y)$: its exceptional fibre at
each point of $C$ is the projective space of normal directions.

Choose two meromorphic sections of $N_{C/U}$ which form a basis at the
generic point. Such sections exist by twisting by a sufficiently positive
line bundle on the curve and applying Riemann--Roch. They give an actual
bundle trivialization away from a finite set, and hence a birational map
\[
                         \PP(N_{C/U})\dashrightarrow C\times\PP^1.
\tag{C.RU.6}
\]
The same meromorphic-frame map applies to every rank-two bundle on $C$.
It therefore identifies this birational class with that of every ruled
surface over this specified curve. In genus zero it also identifies the
rational class, for example through the affine-coordinate birational map
$\PP^2\dashrightarrow\PP^1\times\PP^1$.

We retain the exact Euler changes, including all singularities along the
centres. The reduced inverse image of a point centre replaces that point
by $\PP^2$, adding two to the support Euler number. The reduced inverse
image of the smooth curve $C$ replaces each point by $\PP^1$, adding
$e(C)=2-2g$. These are the proper-map comparisons of (C.ZC.8), which do
not impose a normal-crossing or reduced full-fibre hypothesis. If $p$ is
the total number of point blowups, including the initial one, the support
after blowing up $C$ has
\[
                              e=2p+2-2g.
\tag{C.RU.7}
\]
It is even. There is a smooth open part of the original torus component
disjoint from $R$, since only finitely many proper curves and points have
been affected there. Perform the positive point blocks or negative
point--quartic blocks of C.2.UN.4 in such an open part and in the new
exceptional components they create. Each changes Euler number by exactly
$2$ or $-2$. Choosing the indicated finite number makes (C.RU.7) zero.
All these centres are disjoint from $R$, so the displayed smooth component,
its normal projection to $C$, and its coefficient one are unchanged.
All other fibres remain the original smooth tori. Every blowup has
connected fibres, so the composite map retains connected fibres as well.
Compactness and smoothness follow at each step from the actual point or
smooth-curve blowup of the compact smooth threefold.

For the cohomology comparison use the lines convention in (C.RU.1) and
put $\xi=c_1(\OO_{\PP(N)}(1))$. Pullback and cup product give the integral
projective-bundle map
\[
 H^k(C,\Z)\oplus H^{k-2}(C,\Z)
       \longrightarrow H^k(R,\Z),\qquad
 (a,b)\longmapsto\pi^*a+\xi\smile\pi^*b.
\tag{C.RU.8}
\]
It is an isomorphism: on every fibre the classes $1,\xi$ are the standard
integral basis of $H^*(\PP^1,\Z)$, so the fibrewise basis glues and the
Leray--Hirsch comparison has these exact maps. The relation for this lines
convention is $\xi^2+\pi^*c_1(N)\xi=0$. Since the curve cohomology is
torsion-free, (C.RU.8) proves the integral assertion. The fibrewise
calculation $H^0(\PP^1,\OO)=\C$, $H^1(\PP^1,\OO)=0$ gives
$\pi_*\OO_R=\OO_C$, $R^1\pi_*\OO_R=0$, with their restriction and
evaluation maps. Leray therefore gives the coherent tuple in (C.RU.2).

There is an explicit all-genus computation supporting the construction.
For $g\ge1$ take the plane curve
\[
                 y^2z^{2g-1}=x^{2g+1}-z^{2g+1}.
\tag{C.RU.9}
\]
It is smooth on $z\ne0$. At its single point at infinity, use
$u=x/y,v=z/y$; the equation is $v^{2g-1}(1+v^2)=u^{2g+1}$.
The holomorphic change $v'=v(1+v^2)^{1/(2g-1)}$ identifies its germ with
the plane branch $v'^{2g-1}=u^{2g+1}$. The two exponents are coprime.
Successive point blowups subtract the smaller exponent from the larger,
and the strict curve is smooth when the smaller exponent becomes one.
Summing (C.RU.5) through these actual charts gives total normalization
defect $2g(g-1)$, so its normalization has genus $g$.
The executed code \path{code/root_ruled_occurrence.py} checks all exponent
charts, genus drops, Euler changes, balancing counts and the integral
projective-bundle ranks for $g=1,\ldots,50$. The displayed subtraction
and adjunction formulas prove the same identities for every $g$.
\end{proof}


% C.2.UN.43 from part_C2_secondary_hopf_occurrence.tex

